\documentclass{axiobac}

\annee{1984}
\mois{Juin}
\session{principale}
\epreuve{Mathématiques}
\sectionbac{Math-Sciences \& Math-Technique}
\sections{mathematiques, sciences-techniques}
\duree{4 heures}
\coefficient{5}
\ministere{Ministère de l'Éducation Nationale}

\begin{document}
\entete

%----------------------------------------------------------------------------------------
%   EXERCICE 1
%----------------------------------------------------------------------------------------
\begin{exercice}{4}
\titre{Homothétie et similitude dans le plan complexe}
\chapitres{Similitudes, Nombres complexes}
\banque{0000001}

Soit $\mathcal{P}$ un plan affine euclidien orienté rapporté à un repère orthonormé $(O, \vec{i}, \vec{j})$. \\
On donne les points $A(1,0)$, $B(-1,0)$ et $C(-1,-1)$. \\
À tout point $M$ d'affixe $z$, on fait correspondre :
\begin{itemize}
    \item Le point $M_1$ d'affixe $z_1$ par l'homothétie de centre $A$ et de rapport 2.
    \item Le point $M_2$ d'affixe $z_2$ par la similitude directe de centre $B$, de rapport 2 et d'angle $\frac{\pi}{2}$.
\end{itemize}

\begin{enumerate}
    \item Calculer $z_1$ et $z_2$ en fonction de $z$.
    \item Soit $M'$ le milieu de $[M_1 M_2]$.
    \begin{enumerate}
        \item Calculer l'affixe $z'$ du point $M'$ en fonction de $z$.
        \item Caractériser l'application $T$ qui au point $M$ associe $M'$.
    \end{enumerate}
    \item Démontrer que, pour $M \neq C$, le triangle $CM'M$ est rectangle isocèle.
\end{enumerate}

\end{exercice}

%----------------------------------------------------------------------------------------
%   EXERCICE 2
%----------------------------------------------------------------------------------------
\begin{exercice}{4}
\titre{Application affine, translation et symétrie}
\horsprogramme

$\mathcal{P}$ désigne le plan affine muni d'un repère $(O, \vec{i}, \vec{j})$. \\
Soit $f$ l'application affine de $\mathcal{P}$ qui, au point $M(x, y)$, associe le point $M'(x', y')$ tel que :
\[
\begin{cases}
    x' = x + y \\
    y' = -y + 4
\end{cases}
\]

\begin{enumerate}
    \item Démontrer que $f \circ f$ est une translation de vecteur $\vec{u}$ que l'on déterminera.
    \item Démontrer que $f = t \circ s = s \circ t$ où $t$ est une translation et $s$ une symétrie. Caractériser $t$ et $s$.
\end{enumerate}

\end{exercice}

%----------------------------------------------------------------------------------------
%   PROBLÈME
%----------------------------------------------------------------------------------------
\begin{probleme}{12}

\partie{I}
\titre{Fonction exponentielle, bijection réciproque et aire}
\chapitres{Fonctions exponentielles, Fonctions réciproques, Intégrales}
\banque{0000140}
Soit $f_m$ la fonction numérique à variable réelle définie par :
\[ f_m(x) = \frac{m e^{2x}}{e^{2x} + 1} \]
où $m$ est un paramètre réel de l'intervalle $]0, 2]$.

\begin{enumerate}
    \item 
    \begin{enumerate}
        \item Étudier les variations de $f_m$. En déduire l'ensemble $J_m$ image de $\mathbb{R}$ par $f_m$.
        \item Montrer que $f_m$ est une bijection de $\mathbb{R}$ sur $J_m$.
        \item Calculer $f_m^{-1}(x)$ pour $x \in J_m$.
        \item Quelle est l'image de $\mathbb{R}^*_-$ par $f_m$ ?
    \end{enumerate}

    \item Pour tout $x$ de $\mathbb{R}$, on pose $g_m(x) = f_m(x) - x$.
    \begin{enumerate}
        \item Étudier les variations de $g_m$.
        \item Montrer que l'équation $f_m(x) = x$ admet dans $\mathbb{R}$ une solution unique $x_m$ et que :
        \[ x_m \in \left] \frac{m}{2}, m \right[ \]
    \end{enumerate}

    \item Soit $(\mathcal{C}_m)$ la courbe représentative de $f_m$ dans un plan rapporté à un repère orthonormé $(O, \vec{i}, \vec{j})$ et $(\mathcal{C}'_m)$ celle de $f_m^{-1}$.
    \begin{enumerate}
        \item Étudier la position de $(\mathcal{C}_m)$ et de la droite $\mathcal{D}$ d'équation $y=x$. Construire $(\mathcal{C}_2)$ et $(\mathcal{C'}_2)$.
        \item Soit $(\mathcal{H}_m)$ la partie du plan limitée par $(\mathcal{C}_m)$, $(\mathcal{C}'_m)$ et les droites d'équations respectives $x=0$ et $y=0$. Calculer l'aire $\mathcal{A}$ de $(\mathcal{H}_m)$ en fonction de $x_m$ et $m$.
        \item Montrer que $\lim\limits_{m \to 0^+} \dfrac{x_m}{m} = \dfrac{1}{2}$. En déduire $\lim\limits_{m \to 0^+} \dfrac{\mathcal{A}}{m^2}$.
    \end{enumerate}
\end{enumerate}

\partie{II}
\titre{Endomorphismes, matrices et suite}
\horsprogramme
Soit $E$ un plan vectoriel réel de base $\mathcal{B} = (\vec{i}, \vec{j})$.
À tout couple $(a, t)$ de $\mathbb{R}^2$, on fait correspondre l'endomorphisme $\varphi_{(a,t)}$ de $E$ défini par sa matrice $M_{(a,t)}$ dans la base $\mathcal{B}$ :
\[
M_{(a,t)} = \begin{pmatrix}
    f_2(t) - 2a & 2a \\
    f_2(t) - 3a & 3a
\end{pmatrix}
\]

\begin{enumerate}
    \item 
    \begin{enumerate}
        \item Déterminer l'ensemble $S$ des couples de réels $(a, t)$ pour chacun desquels $\varphi_{(a,t)}$ est un automorphisme de $E$.
        \item Dans le cas où $\varphi_{(a,t)}$ n'est pas un automorphisme de $E$, déterminer le noyau et l'image de $\varphi_{(a,t)}$. Montrer dans ce cas que $\varphi_{(a,t)}$ est la composée de deux applications simples que l'on caractérisera.
    \end{enumerate}

    \item Dans toute la suite, on prend $a=1$, $t \in \mathbb{R}^*_-$ et on note $T = f_2(t)$ et $\varphi_T = \varphi_{(1,t)}$. \\
    On pose $\varphi_T^1 = \varphi_T$ et pour tout $n \in \mathbb{N}^*$, $\varphi_T^{n+1} = \varphi_T^n \circ \varphi_T$.
    \begin{enumerate}
        \item Démontrer, par récurrence, que pour tout $n$ de $\mathbb{N}^*$, la matrice dans la base $\mathcal{B}$ de $\varphi_T^n$ est de la forme :
        \[
        \begin{pmatrix}
            T^n - P_n(T) & P_n(T) \\
            T^n - P_n(T) - 1 & P_n(T) + 1
        \end{pmatrix}
        \]
        où $P_n(T)$ est un polynôme vérifiant pour tout $n$ de $\mathbb{N}^*$ :
        \[ P_{n+1}(T) = P_n(T) + 2T^n \quad (1) \]
        \item Utiliser la relation (1) pour expliciter le polynôme $P_n(T)$ pour tout $n$ de $\mathbb{N}^*$.
    \end{enumerate}

    \item 
    \begin{enumerate}
        \item Montrer que la suite de terme général $P_n(T)$ converge et calculer sa limite $l(T)$.
        \item Calculer l'intégrale $\int_{-\ln 4}^{-\ln 2} l(f_2(t)) \, dt$.
    \end{enumerate}
\end{enumerate}

\end{probleme}

\end{document}
